%% EXEMPLO FICTÍCIO — Atividade 2: simulação computacional
\begin{atividade}{Simulação do estágio de saída e do limite de tensão de conformidade}{
  tipo         = simulacao,
  modalidade   = PES,
  alternativa  = VPC,
  horas        = 45,
  inicio       = 2024-02,
  fim          = 2024-04,
  projeto      = {NEBULA, estimulação elétrica funcional para ciclismo},
  papel        = Responsável,
  supervisor   = {Rafael Fictício, supervisor do projeto NEBULA},
  registros    = {NBL-109; NRO-PRO-013-901},
  competencias = {II, III, VIII},
  disciplinas  = {Circuitos Elétricos II; Eletrônica Analógica; Sistemas de Controle; Cálculo Numérico},
}

\begin{contexto}
Antes de refazer o estágio de saída (Atividade~1), a equipe precisava responder a duas perguntas.
Primeira: com a impedância real dos eletrodos, que tensão de alimentação garante
\SI{100}{\milli\ampere} durante um pulso inteiro de \SI{300}{\micro\second}? Segunda: o laço do
amplificador operacional com o MOSFET, que oscilava na revisão A, fica estável com os componentes
novos? Da primeira dependia a escolha entre manter \SI{90}{\volt} e subir para \SI{120}{\volt} ---
com MOSFETs, capacitores e isolação mais caros; da segunda, o capacitor de compensação.
\end{contexto}

\begin{contribuicao}
Montei o modelo da carga a partir da literatura, levantei os parâmetros dos componentes, rodei as
simulações e as varreduras, validei o resultado contra a medição na placa e escrevi o registro de
decisão de arquitetura que fixou a tensão em \SI{120}{\volt} (\texttt{NRO-PRO-013-901}). O modelo do
MOSFET é o do fabricante; o do amplificador operacional, o macromodelo do fabricante, sem alteração.
\end{contribuicao}

\begin{desenvolvimento}
\fichatecnica{
  ferramenta    = {LTspice 24},
  fenomeno      = {A tensão no eletrodo durante o pulso e a resposta da fonte de corrente à borda do pulso},
  modelo        = {Circuito elétrico: fonte de corrente da placa e carga no modelo de Randles; análise transitória e de pequenos sinais},
  parametros    = {$R_s = 200\,\Omega$; $R_p$ de 500\,$\Omega$ a 2\,k$\Omega$; $C_p = 100$\,nF (faixas de Merrill, Bikson e Jefferys, 2005); modelos dos fabricantes para U1 e Q1},
  discretizacao = {Passo máximo de 10\,ns; com 5\,ns, o sobressinal mudou menos de 0,2 ponto percentual},
  validacao     = {Comparação com a placa montada, em carga resistiva de 1\,k$\Omega$ (Atividade~4)},
  arquivos      = {Repositório \texttt{nebula-sim}, pasta \texttt{saida/}},
}

\textbf{O modelo da carga.} O conjunto eletrodo e pele se comporta, em primeira aproximação, como
um resistor em série com um paralelo RC --- o modelo de Randles (Figura~\ref{fig:ex-randles}) ---, com
valores na faixa relatada por \textcite{merrill2005}. Sob um degrau de corrente $I$, a tensão no
eletrodo é
\begin{equation}
  v_e(t) = I \left[ R_s + R_p \left( 1 - e^{-t/(R_p C_p)} \right) \right].
  \label{eq:ex-randles}
\end{equation}
A fonte mantém a corrente enquanto $v_e$ somada às quedas em Q1 e em $R_S$ (cerca de
\SI{2}{\volt}) não passa da tensão de alimentação: com \SI{120}{\volt}, a conformidade é de
$V_c \approx \SI{118}{\volt}$. Resolvendo \eqref{eq:ex-randles} para $v_e = V_c$, o instante em que a
corrente começa a cair é
\begin{equation}
  t_\mathrm{sat} = -R_p C_p \ln\!\left( 1 - \frac{V_c/I - R_s}{R_p} \right),
  \label{eq:ex-tsat}
\end{equation}
que só existe quando $I (R_s + R_p) > V_c$.

\begin{figure}[htbp]
  \centering
  % O modelo da carga (exemplo fictício): a fonte de corrente sobre o
% conjunto eletrodo e pele, no modelo de Randles.
\begin{circuitikz}[font=\footnotesize, line width=0.5pt, scale=0.85, transform shape]
  \draw (0,0) to[I, l=$i(t)$] (0,3) to[R, l=$R_s$, -*] (3,3) coordinate (a);
  \draw (a) -- (3,3.9) to[R, l=$R_p$] (6.2,3.9) -- (6.2,3) coordinate (b);
  \draw (a) -- (3,2.1) to[C, l_=$C_p$] (6.2,2.1) -- (b);
  \draw (b) to[short, *-] (7.6,3) -- (7.6,0) -- (0,0);
  \draw (7.6,0) node[ground] {};
  \draw[nro-cinza, densely dashed, rounded corners=2pt] (1.3,1.5) rectangle (6.9,4.8);
  \node[nro-cinza, font=\scriptsize, anchor=south] at (4.1,4.8) {eletrodo e pele (modelo de Randles)};
  \draw[-{Stealth[length=1.6mm]}, nro-axon] (8.4,0.3) -- node[right, text=nro-axon] {$v_e(t)$} (8.4,2.9);
\end{circuitikz}

  \caption{O modelo da carga: a fonte de corrente sobre o eletrodo e a pele (modelo de Randles).}
  \fonte{Elaborado pela autora, a partir de \textcite{merrill2005}.}
  \label{fig:ex-randles}
\end{figure}

\textbf{A estabilidade.} Na análise de pequenos sinais, o laço de U1 com Q1 tinha
\SI{38}{\degree} de margem de fase: a capacitância de entrada do MOSFET novo, maior que a do antigo,
somada à resistência de gate, punha um polo perto do cruzamento. Com $C_c = \SI{22}{\pico\farad}$
entre a saída e a entrada inversora de U1, a margem subiu para \SI{62}{\degree}. Testei 10, 22 e
\SI{47}{\pico\farad}: \SI{47}{\pico\farad} amortecia mais, mas alongava a subida além do que o pulso
de \SI{50}{\micro\second} tolera.
\end{desenvolvimento}

\begin{resultados}
A Figura~\ref{fig:ex-tensao} responde à primeira pergunta. Com \SI{120}{\volt}, o pulso de
\SI{300}{\micro\second} a \SI{100}{\milli\ampere} se sustenta até $R_p = \SI{1}{\kilo\ohm}$
($v_e = \SI{115}{\volt}$ no fim do pulso); com $R_p = \SI{2}{\kilo\ohm}$, a corrente cai a partir de
\SI{135}{\micro\second}, como prevê \eqref{eq:ex-tsat}. Com \SI{90}{\volt}, o limite seria
$R_p \approx \SI{690}{\ohm}$ --- abaixo dos eletrodos já gastos que a equipe usa. A decisão foi subir
para \SI{120}{\volt} e, para impedâncias maiores, fazer o firmware medir a impedância no começo da
sessão e limitar a amplitude: isso virou requisito na revisão C da USRS.

\begin{figure}[htbp]
  \centering
  % Tensão no eletrodo durante um pulso de 100 mA, pela eq. (2), limitada
% pela tensão de conformidade (exemplo fictício).
\begin{tikzpicture}
  \begin{axis}[
      width=0.86\linewidth, height=60mm,
      xlabel={Tempo desde o início do pulso (\si{\micro\second})},
      ylabel={Tensão no eletrodo (\si{\volt})},
      xmin=0, xmax=300, ymin=0, ymax=135,
      grid=major, grid style={nro-fundo},
      tick label style={font=\scriptsize}, label style={font=\footnotesize},
      legend pos=south east, legend cell align=left,
      legend style={font=\scriptsize, draw=nro-linha, fill=white, fill opacity=0.9, text opacity=1}]
    \addplot[domain=0:300, samples=121, very thick, nro-axon] {min(118, 0.1*(200+500*(1-exp(-x/50))))};
    \addlegendentry{$R_p = 500\,\Omega$}
    \addplot[domain=0:300, samples=121, very thick, nro-tinta, dashed] {min(118, 0.1*(200+1000*(1-exp(-x/100))))};
    \addlegendentry{$R_p = 1\,\mathrm{k}\Omega$}
    \addplot[domain=0:300, samples=121, very thick, nro-pulso-texto] {min(118, 0.1*(200+2000*(1-exp(-x/200))))};
    \addlegendentry{$R_p = 2\,\mathrm{k}\Omega$}
    \addplot[domain=0:300, thin, nro-cinza, densely dotted] {118};
    \addlegendentry{conformidade (118\,V)}
    \draw[nro-pulso-texto] (axis cs:134.7,118) circle (1.6pt);
    \node[font=\scriptsize, text=nro-pulso-texto, anchor=south east] at (axis cs:132,119) {a corrente cai a partir de \SI{135}{\micro\second}};
  \end{axis}
\end{tikzpicture}

  \caption{Tensão no eletrodo durante um pulso de \SI{100}{\milli\ampere}, para três valores de $R_p$,
    com a alimentação de \SI{120}{\volt}.}
  \fonte{Elaborado pela autora (simulação, eq.~\ref{eq:ex-randles}).}
  \label{fig:ex-tensao}
\end{figure}

A Figura~\ref{fig:ex-borda} responde à segunda. Sem $C_c$, a simulação mostra 18\,\% de sobressinal
--- a oscilação que a revisão A tinha; com \SI{22}{\pico\farad}, 3\,\%. Na placa montada, o sobressinal
medido foi de 1,9\,\% e o tempo de subida (10 a 90\,\%), de \SI{1,8}{\micro\second}, contra
\SI{1,7}{\micro\second} na simulação: 6\,\% de diferença. A amplitude simulada em \SI{1}{\kilo\ohm}
(\SI{100,0}{\milli\ampere}) ficou a 0,9\,\% da medida.

\begin{figure}[htbp]
  \centering
  % A borda de subida da corrente: simulação sem e com a compensação, e a
% medição na placa montada (exemplo fictício).
\begin{tikzpicture}
  \begin{axis}[
      width=0.86\linewidth, height=58mm,
      xlabel={Tempo (\si{\micro\second})},
      ylabel={Corrente (\% da programada)},
      xmin=0, xmax=8, ymin=0, ymax=125,
      grid=major, grid style={nro-fundo},
      tick label style={font=\scriptsize}, label style={font=\footnotesize},
      legend pos=south east, legend cell align=left,
      legend style={font=\scriptsize, draw=nro-linha, fill=white, fill opacity=0.9, text opacity=1}]
    \addplot[domain=0:8, samples=161, thick, nro-pulso-texto]
      {100*(1-exp(-0.48*1.35*x)*(cos(deg(1.35*sqrt(1-0.48^2)*x)) + 0.48/sqrt(1-0.48^2)*sin(deg(1.35*sqrt(1-0.48^2)*x))))};
    \addlegendentry{simulada, sem $C_c$ (sobressinal de 18\,\%)}
    \addplot[domain=0:8, samples=161, very thick, nro-axon]
      {100*(1-exp(-0.74*1.35*x)*(cos(deg(1.35*sqrt(1-0.74^2)*x)) + 0.74/sqrt(1-0.74^2)*sin(deg(1.35*sqrt(1-0.74^2)*x))))};
    \addlegendentry{simulada, $C_c = 22$\,pF (3\,\%)}
    \addplot[only marks, mark=*, mark size=1.4pt, nro-tinta]
      coordinates {(0.5,13.8) (1,40.2) (1.5,66.6) (2,83.5) (2.5,95.7) (3,100.8) (4,101.9) (5,101.4) (6,99.2) (7,99.8) (8,98.9)};
    \addlegendentry{medida na placa, $C_c = 22$\,pF}
  \end{axis}
\end{tikzpicture}

  \caption{A borda de subida da corrente: simulada sem e com a compensação, e medida na placa.}
  \fonte{Elaborado pela autora.}
  \label{fig:ex-borda}
\end{figure}

\textbf{As limitações.} O modelo de Randles é linear, e a impedância real do eletrodo muda com a
densidade de corrente, com o suor e com o desgaste do gel; os parâmetros vieram da literatura e de
medições em placa de teste, não no usuário. A simulação serve para dimensionar a alimentação e a
compensação; não permite concluir nada sobre o conforto nem sobre a segurança da pele, que ficam com o
protocolo clínico.
\end{resultados}

\begin{evidencias}
\evidencia{Repositório}{nebula-sim, pasta saida/}{Os arquivos da simulação (.asc), os parâmetros e as varreduras}{GitHub da equipe (acesso a pedido)}
\evidencia{Arquivo controlado}{NRO-PRO-013-901}{O registro da decisão de subir a alimentação para 120\,V, com as alternativas}{SOMA › Arquivos}
\evidencia{Atividade}{NBL-109}{O cartão da simulação, atribuído a mim e concluído}{SOMA › Atividades › NEBULA}
\end{evidencias}

\begin{aprendizados}
A simulação juntou três disciplinas numa pergunta só: o transitório de primeira ordem de Circuitos
Elétricos II está em \eqref{eq:ex-randles}, a margem de fase de Sistemas de Controle decidiu o
capacitor, e o cuidado com o passo de tempo veio de Cálculo Numérico. O que aprendi fora da sala: um
modelo vale o que valem os seus parâmetros, e a simulação só ganhou crédito na equipe quando bateu com
a medição. Hoje eu faria a validação antes de apresentar o resultado, e não depois.
\end{aprendizados}
\end{atividade}
